% GENERATED FILE. Edit canonical lessons, then run npm run generate:books.
% canonical-source-hash: fnv1a64:acb6883de595630d
% canonical-lessons: ZH-C184-dujia, ZH-C184-jueding, ZH-C184-canguan, ZH-C184-laojia, ZH-C184-fangsong

\chapter{Holiday Plans}
\label{ch:zh-holiday-plans}
\hlchaptermodality{184}

\begin{chapteropening}
\textbf{By the end of this chapter:} \emph{I can talk about going on holiday, deciding where, looking round places, my home town and relaxing.}

Complete the last lesson of chapter 184: I can talk about going on holiday, deciding where, looking round places, my home town and relaxing.
\end{chapteropening}

\section[dùjià]{dùjià — to go on holiday; to spend a holiday}

\label{lesson:ZH-C184-dujia}



\begin{quote}
Before the new one: say the Mandarin for to lock up, then the Mandarin for to lie down.
\end{quote}

\subsection*{You'll want to know: dùjià}
\emph{dùjià} — \textquotedblleft{}to go on holiday; to spend a holiday\textquotedblright{}.

\begin{quote}
\emph{xià ge yuè wǒmen qù Běijīng dùjià} — Next month we're going on holiday to Beijing.
\end{quote}

Its characters come later; for now it is a word for the ear.

\begin{cousinweb}
\emph{dùjià} is built from \emph{dù} and \emph{jià} 'leave', both inherited from Middle Chinese, \emph{duH} and \emph{kaeH} (Guangyun, in Baxter's transcription). The dictionary glosses the second 'leave from duty' and the first 'a rule'; 'to spend time' is not given there.
\end{cousinweb}

\begin{grammarlens}[title={What you've built}]
One more word for talking about plans.
\end{grammarlens}

\subsection*{Your turn}
\begin{itemize}
\raggedright
  \item \emph{Say it:} \emph{dùjià}
  \item \emph{Say it:} \emph{dùjià}, once more
  \item \emph{Say it:} \emph{xià ge yuè wǒmen qù Běijīng dùjià} — Next month we're going on holiday to Beijing.
  \item \emph{Recall:} say the Mandarin for to lock up, then the Mandarin for to lie down, then say \emph{dùjià} again
\end{itemize}

\begin{tcolorbox}[breakable,colback=teal!4,colframe=teal!35!black,title={Before you move on}]
What does dùjià mean? (\textquotedblleft{}To go on holiday; to spend a holiday\textquotedblright{}.) Say it once more.
\end{tcolorbox}
\section[juédìng]{juédìng — to decide; a decision}

\label{lesson:ZH-C184-jueding}



\begin{quote}
Before the new one: say the Mandarin for to lie down, then the Mandarin for to go on holiday.
\end{quote}

\subsection*{You'll want to know: juédìng}
\emph{juédìng} — \textquotedblleft{}to decide; a decision\textquotedblright{}.

\begin{quote}
\emph{wǒmen hái méi juédìng qù nǎr} — We haven't decided where to go yet.
\end{quote}

Its characters come later; for now it is a word for the ear.

\begin{cousinweb}
\emph{juédìng} is built from \emph{jué} and \emph{dìng}, both inherited from Middle Chinese, \emph{kwet} and \emph{dengH} (Guangyun, in Baxter's transcription). The dictionary glosses the first 'to cut off' and 'to break', and the second 'settled': to cut off and settle.
\end{cousinweb}

\begin{grammarlens}[title={What you've built}]
One more word for talking about plans.
\end{grammarlens}

\subsection*{Your turn}
\begin{itemize}
\raggedright
  \item \emph{Say it:} \emph{juédìng}
  \item \emph{Say it:} \emph{juédìng}, once more
  \item \emph{Say it:} \emph{wǒmen hái méi juédìng qù nǎr} — We haven't decided where to go yet.
  \item \emph{Recall:} say the Mandarin for to lie down, then the Mandarin for to go on holiday, then say \emph{juédìng} again
\end{itemize}

\begin{tcolorbox}[breakable,colback=teal!4,colframe=teal!35!black,title={Before you move on}]
What does juédìng mean? (\textquotedblleft{}To decide; a decision\textquotedblright{}.) Say it once more.
\end{tcolorbox}
\section[cānguān]{cānguān — to look round a place; to visit (a sight)}

\label{lesson:ZH-C184-canguan}



\begin{quote}
Before the new one: say the Mandarin for to go on holiday, then the Mandarin for to decide.
\end{quote}

\subsection*{You'll want to know: cānguān}
\emph{cānguān} — \textquotedblleft{}to look round a place; to visit (a sight)\textquotedblright{}.

\begin{quote}
\emph{míngtiān wǒmen qù cānguān dàxué} — Tomorrow we're going to look round the university.
\end{quote}

Its characters come later; for now it is a word for the ear.

\begin{cousinweb}
\emph{cānguān} is built from \emph{cān} 'to present oneself' and \emph{guān} 'to look', both inherited from Middle Chinese, \emph{tshom} and \emph{kwan} (Guangyun, in Baxter's transcription).
\end{cousinweb}

\begin{grammarlens}[title={What you've built}]
One more word for talking about plans.
\end{grammarlens}

\subsection*{Your turn}
\begin{itemize}
\raggedright
  \item \emph{Say it:} \emph{cānguān}
  \item \emph{Say it:} \emph{cānguān}, once more
  \item \emph{Say it:} \emph{míngtiān wǒmen qù cānguān dàxué} — Tomorrow we're going to look round the university.
  \item \emph{Recall:} say the Mandarin for to go on holiday, then the Mandarin for to decide, then say \emph{cānguān} again
\end{itemize}

\begin{tcolorbox}[breakable,colback=teal!4,colframe=teal!35!black,title={Before you move on}]
What does cānguān mean? (\textquotedblleft{}To look round a place; to visit (a sight)\textquotedblright{}.) Say it once more.
\end{tcolorbox}
\section[lǎojiā]{lǎojiā — home town; family home}

\label{lesson:ZH-C184-laojia}



\begin{quote}
Before the new one: say the Mandarin for to decide, then the Mandarin for to look round a place.
\end{quote}

\subsection*{You'll want to know: lǎojiā}
\emph{lǎojiā} — \textquotedblleft{}home town; family home\textquotedblright{}.

\begin{quote}
\emph{Chūnjié wǒ yào huí lǎojiā} — I'm going back to my home town for the new year holiday.
\end{quote}

Its characters come later; for now it is a word for the ear.

\begin{cousinweb}
\emph{lǎojiā} is built from \emph{lǎo} 'old' and \zh{家} 'home', both inherited from Middle Chinese, \emph{lawX} and \emph{kae} (Guangyun, in Baxter's transcription): the old home.
\end{cousinweb}

\begin{grammarlens}[title={What you've built}]
One more word for talking about plans.
\end{grammarlens}

\subsection*{Your turn}
\begin{itemize}
\raggedright
  \item \emph{Say it:} \emph{lǎojiā}
  \item \emph{Say it:} \emph{lǎojiā}, once more
  \item \emph{Say it:} \emph{Chūnjié wǒ yào huí lǎojiā} — I'm going back to my home town for the new year holiday.
  \item \emph{Recall:} say the Mandarin for to decide, then the Mandarin for to look round a place, then say \emph{lǎojiā} again
\end{itemize}

\begin{tcolorbox}[breakable,colback=teal!4,colframe=teal!35!black,title={Before you move on}]
What does lǎojiā mean? (\textquotedblleft{}Home town; family home\textquotedblright{}.) Say it once more.
\end{tcolorbox}
\section[fàngsōng]{fàngsōng — to relax; to loosen}

\label{lesson:ZH-C184-fangsong}



\begin{quote}
Before the new one: say the Mandarin for to look round a place, then the Mandarin for one's home town.
\end{quote}

\subsection*{You'll want to know: fàngsōng}
\emph{fàngsōng} — \textquotedblleft{}to relax; to loosen\textquotedblright{}.

\begin{quote}
\emph{wǒ xiǎng fàngsōng bù xiǎng chūqù} — I want to relax; I don't feel like going out.
\end{quote}

Its characters come later; for now it is a word for the ear.

\begin{cousinweb}
\emph{fàngsōng} is built from \emph{fàng} 'to let go' and \emph{sōng} 'loose', both inherited from Middle Chinese, \emph{pjangH} and \emph{sowng} (Guangyun, in Baxter's transcription), which glosses the second as hair hanging loose.
\end{cousinweb}

\begin{grammarlens}[title={What you've built}]
One more word for talking about plans.
\end{grammarlens}

\subsection*{Your turn}
\begin{itemize}
\raggedright
  \item \emph{Say it:} \emph{fàngsōng}
  \item \emph{Say it:} \emph{fàngsōng}, once more
  \item \emph{Say it:} \emph{wǒ xiǎng fàngsōng bù xiǎng chūqù} — I want to relax; I don't feel like going out.
  \item \emph{Recall:} say the Mandarin for to look round a place, then the Mandarin for one's home town, then say \emph{fàngsōng} again
\end{itemize}

\begin{tcolorbox}[breakable,colback=teal!4,colframe=teal!35!black,title={Before you move on}]
What does fàngsōng mean? (\textquotedblleft{}To relax; to loosen\textquotedblright{}.) Say it once more.
\end{tcolorbox}
